% =============================================================================
% MLSys·im Tutorial — Module 4: Design Space Exploration & Synthesis
% =============================================================================
\documentclass[aspectratio=169, 12pt]{beamer}
\usepackage{../../../slides/assets/beamerthememlsys}

\mlsyssetup{
  volume       = {Tutorial},
  chapter      = {Module 4},
  logo         = {../../../slides/assets/img/logo-mlsysbook.png},
  instlogo     = {../../../slides/assets/img/logo-harvard.png},
  chaptertitle = {MLSys·im: DSE \& Synthesis},
}

% --- Fonts & Packages ---
\usepackage[T1]{fontenc}
\usepackage[scaled=0.9]{helvet}
\usepackage{courier}
\renewcommand{\familydefault}{\sfdefault}
\usepackage{amsmath}
\usepackage{booktabs}
\usepackage[table]{xcolor}
\usepackage{listings}
\usepackage{tikz}

% --- Code listings ---
\lstset{
  language=Python,
  basicstyle=\ttfamily\footnotesize,
  keywordstyle=\color{crimson}\bfseries,
  stringstyle=\color{datastroke},
  commentstyle=\color{midgray}\itshape,
  backgroundcolor=\color{computeblue!20},
  frame=single,
  rulecolor=\color{computestroke},
  numbers=none,
  breaklines=true,
  columns=fullflexible,
  keepspaces=true,
  showstringspaces=false,
  xleftmargin=4pt,
  xrightmargin=4pt,
  aboveskip=6pt,
  belowskip=4pt,
}

\newcommand{\mlsysim}{\texttt{mlsysim}}
\graphicspath{{images/}}

\title{MLSys·im Tutorial --- Module 4}
\subtitle{Design Space Exploration \& Synthesis}
\author{Vijay Janapa Reddi}
\institute{Harvard University}
\date{Conference Tutorial}

\begin{document}

\begin{frame}[plain]
  \titlepage
\end{frame}

\begin{frame}{Roadmap: Conference Tutorial}
\centering\small
\begin{tabular}{rll}
\toprule
\textbf{Module} & \textbf{Topic} & \textbf{Status} \\
\midrule
Module 1 & Foundations \& Architecture & \checkmark Done \\
Module 2 & Advanced Single-Node Analysis & \checkmark Done \\
Module 3 & Scale, Dollars, and Carbon & \checkmark Done \\
\rowcolor{crimson!12}
\textbf{Module 4} & \textbf{Design Space Exploration \& Synthesis} & \textbf{$\leftarrow$ You are here} \\
\bottomrule
\end{tabular}
\end{frame}

\section{Rapid Parametric Sweeps}

\begin{frame}{The Combinatorial Explosion}
Finding the optimal serving configuration requires testing:
\[
|\text{hardware}| \times |\text{batch sizes}| \times |\text{precisions}| \times |\text{parallelism configs}|
\]
This space easily exceeds $10^4$ configurations. Because \mlsysim{} uses analytical math (not cycle-accurate simulation), each evaluation takes $<1$\,ms. 
\end{frame}

\begin{frame}[fragile]{Live Demo: Programmatic Sweeps (Scenario A)}
\begin{lstlisting}[language=Python]
from mlsysim.engine.engine import Engine
from mlsysim.hardware.registry import Hardware
from mlsysim.models.registry import Models
import pandas as pd

model = Models.Language.Llama3_8B
hardware = Hardware.Cloud.H100

results = []
for batch_size in [1, 8, 32, 128, 256]:
    perf = Engine.solve(model, hardware, batch_size=batch_size, precision="fp16")
    results.append({
        "Batch Size": batch_size,
        "Throughput (tok/s)": perf.throughput.magnitude,
        "Bottleneck": perf.bottleneck
    })

print(pd.DataFrame(results))
\end{lstlisting}
\end{frame}

\section{TinyML to Frontier}

\begin{frame}{Same Roofline, 9 Orders of Magnitude}
\begin{columns}[T]
\column{0.52\textwidth}
\centering
\includegraphics[width=0.85\textwidth]{images/pdf/hardware-spectrum.pdf}

\column{0.45\textwidth}
\scriptsize
\begin{tabular}{lrr}
\toprule
\textbf{Device} & \textbf{FLOPS} & \textbf{TDP} \\
\midrule
nRF52840  & 64\,M  & 15\,mW \\
ESP32-S3  & 500\,M & 400\,mW \\
\rowcolor{gray!15}
H100 SXM  & 989\,T & 700\,W \\
\bottomrule
\end{tabular}

\vspace{0.5em}
\begin{tabular}{lr}
\textbf{Compute Range} & $\sim 10^{7}\times$ \\
\textbf{Power Range}   & $\sim 10^{4.7}\times$ \\
\end{tabular}
\end{columns}
\vfill
\begin{center}
\alert{The Roofline Model is universal. The physics apply identically to a \$2 Microcontroller and a \$3M GPU Rack.}
\end{center}
\end{frame}

\section{Sensitivity Analysis}

\begin{frame}[fragile]{Wall 21: Sensitivity Analysis}
\note{[3 min] ``Which knob should I turn next?'' The parameter with the largest partial derivative.}
\begin{lstlisting}
from mlsysim.engine.solver import SensitivitySolver
solver = SensitivitySolver()

result = solver.solve(
    model=Models.Language.Llama3_8B, hardware=Hardware.Cloud.H100, 
    precision="fp16", efficiency=0.5)

print(f"Binding Constraint: {result.binding_constraint}")
for param, sensitivity in result.sensitivities.items():
    tag = "<<<" if param == result.binding_constraint else ""
    print(f"  {param:>20}: {sensitivity:+.4f}  {tag}")
\end{lstlisting}

\vspace{0.3em}
\small
\textbf{The Golden Rule:} Invest in the parameter with the \emph{largest} partial derivative. Improving a non-binding parameter yields \textbf{zero} measurable gain.
\end{frame}

\section{SLA-Driven Synthesis}

\begin{frame}[fragile]{Live Demo: Inverting the Roofline (Scenario D)}
Instead of asking "How fast is this GPU?", what if we ask "What hardware do I need to buy to meet my 30ms latency SLA?"
\begin{lstlisting}[language=Python]
from mlsysim.solvers import SynthesisSolver
from mlsysim.models.registry import Models
from mlsysim.core.constants import Q_

solver = SynthesisSolver()
requirements = solver.solve(
    model=Models.Language.Llama3_8B,
    target_latency=Q_("30 ms"),
    batch_size=1,
    precision="fp16"
)

print(f"Required HBM Bandwidth: {requirements.required_bw.to('GB/s'):.1f}")
print(f"Required Compute:       {requirements.required_flops.to('TFLOPs/s'):.1f}")
\end{lstlisting}
\end{frame}

\section{Capstone \& Wrap-Up}

\begin{frame}[fragile]{Design Challenge: The Capstone}
\begin{alertblock}{The Problem}
\textbf{\$5M budget.} Serve Llama-3 70B at \textbf{1{,}000 QPS} with
\textbf{$<$100\,ms TTFT} in \textbf{two regions} (US-East + EU-West).
Design the fleet.
\end{alertblock}

\vspace{0.3em}
\textbf{You must specify using \mlsysim{}:}
\begin{enumerate}\footnotesize
  \item \textbf{Hardware choice:} Which GPU? How many?
  \item \textbf{Parallelism strategy:} TP $\times$ PP?
  \item \textbf{Precision:} FP16? FP8? INT4?
  \item \textbf{Geographic placement:} Carbon impact?
\end{enumerate}
\end{frame}

\begin{frame}{Resources \& Next Steps}
\begin{columns}[T]
\column{0.55\textwidth}
\textbf{Get Started}
\begin{itemize}
  \item \texttt{pip install mlsysim}
  \item GitHub: \texttt{harvard-edge/mlsysim}
  \item \textbf{Code Cookbook:} See the 5 interactive scenarios in the Appendix!
  \item Full docs: \texttt{mlsysim.readthedocs.io}
\end{itemize}

\vspace{0.5em}
\textbf{The Textbook}
\begin{itemize}
  \item \emph{Machine Learning Systems}
  \item Volume I: Foundations (single node)
  \item Volume II: Systems at Scale (fleet)
  \item \texttt{mlsysbook.ai}
\end{itemize}

\column{0.40\textwidth}
\textbf{Key Papers}
\begin{itemize}
  \item Williams et al.\ (2009)\\
        {\footnotesize Roofline Model}
  \item Chowdhery et al.\ (2022)\\
        {\footnotesize PaLM / MFU}
  \item Hoffmann et al.\ (2022)\\
        {\footnotesize Chinchilla Scaling}
  \item Patterson et al.\ (2021)\\
        {\footnotesize Carbon \& Training}
  \item OpenAI (2024)\\
        {\footnotesize o1 / Reasoning Scaling}
\end{itemize}
\end{columns}

\vspace{1em}
\begin{center}
\Large\bfseries Thank you! Questions?
\end{center}
\end{frame}

\end{document}